\section{Обозначения}
\label{sec:Chapter1} \index{Chapter1}

Под $\Z_m$ будем подразумевать кольцо $\Z/m\Z$ с обычными представителями $\{0, \ldots, m-1\}$. Ими же зададим стандартное вложение $\Z_m \hookrightarrow \Z$, сравнение из которого будем использовать, говоря $x < y$ для $x, y \in \Z_m$.

Пусть $n \in \N$. Обозначим за $\mathcal{R}^n$ кольцо многочленов $\Z[x]/(x^n+1)$. Вспомним, что если $n$ --- степень $2$, то $x^n+1$ будет многочленом деления круга $\reqv$ неприводимым над $\Z$. Если дополнительно число $q$ --- простое, примем за $\mathcal{R}_q^n$ кольцо $\Z_q[x]/(x^n+1)$.

Элементы $\mathcal{R}_q^n$ можно считать многочленами степени $n-1$ с $n$~коэффициентами из $\Z_q$. Будем писать их жирным шрифтом~$\bold a$. При этом их $i$-ый коэффициент будем обозначать $a_i$. Если все коэффициенты независимо генерируются из распределения $P$ или равномерно из мн-ва $S$, будем обозначать это $\bold{a} \leftarrow P$ или $\bold{a} \leftarrow S$ соответственно.

Сумма и произведение естественным образом наследуются в $\mathcal{R}_q^n$ из $\Z_q[x]$, причём из-за выполнения $x^n = -1$ получаем равенство:
\begin{equation}
	(\bold a\bold b)_k = \sum_{i=0}^k{a_ib_{k-i}}\,-\sum_{i=k+1}^{n-1}{a_ib_{n+k-i}}
\end{equation}

Для распределения $P$ будем обозначать вероятность исхода $x$, как $P[x]$. Если при этом $P$ принимает значения в $\{0, \ldots, d-1\}$, ему естественным образом сопоставляется многочлен:
\begin{equation}
	\hat P(z) = \sum_{i=0}^{d-1}{P[i]z^i}
\end{equation}

Для события $A$ примем за $\Pr(A)$ его вероятность, а за $I\{A\}$ --- индикатор его выполнения.

Станем обозначать биномиальное распределение c $n$ попытками и вероятностью успеха $p$, как $\mathcal{B}(n, p)$; оно принимает значения $0 \leq x \leq n$. Если $\mu, \nu \sim \mathcal{B}(k, \frac12)$ --- независимые случайные величины, распределение разности $\mu-\nu$ обозначим за $\chi_k$ ({\it centered binomial distribution}), соответственно, его значения будут $-k \leq x \leq k$.

Свёртку в $\Z_q$ распределений $P$ и $R$ будем писать, как $P * R$. Свёртку же распределения $P$ с собой $n-1$ раз изобразим $P^{\circledast n}$.

Под алфавитом размера $Q$ будем подразумевать $\Sigma_Q = \{0, \ldots, Q-1\}$. При арифметических операциях с его символами будем считать $\Sigma_Q \hookrightarrow \Z_Q$. Введём на $\Sigma_Q$ норму Хэмминга $|s|_H$ и норму Ли $|s|_L$, как:
\begin{equation}
	|s|_H = I\{s \neq 0\}; \qquad |s|_L = \min(s, Q-s)
\end{equation}

После этого естественным образом определим норму на $\Sigma_Q^n$:
\begin{equation}
	\|m\| = \sum_{i=0}^{n-1}{|m_i|}
\end{equation}

Решётками в $\R^n$ станем называть мн-ва векторов $L$, состоящие из целочисленных ЛК ({\it линейных комбинаций}) фиксированных $n$ ЛНЗ ({\it линейно независимых}) векторов --- базиса решётки.
